\section{Introduction}
\label{sec:intro}

\begin{figure}[h]
\centering
\includegraphics[width=\textwidth]{signature_method}
\caption{\textbf{\ours{} at a glance.}
\textbf{(A)} The pipeline: generate nine short stories per emotion class with
the LLM itself, embed them at the LLM's final layer, average per class to
form nine centroids, and take the first PCA component of the mean-centered
centroids as the valence axis $v$.
\textbf{(B)} Zero-shot SST-2 AUC of the V-axis extracted at each transformer
layer of Qwen2.5-1.5B: strong at L1 (embedding-level signal), falls to chance
through layers 2--26 (the ``dilution zone''), and \emph{re-crystallizes
abruptly} at layers 27--28.
\textbf{(C)} Data efficiency: a single story per class already beats a
supervised logistic regression trained on ten labelled SST-2 examples,
and nine stories match a supervised classifier trained on roughly a hundred
labels --- a 10--30$\times$ reduction in labeled data.}
\label{fig:hero}
\end{figure}

Practical affect classification today still depends on labelled data:
thousands of annotations to train a sentiment model, hand-curated
multilingual datasets to support non-English deployment, bespoke contrastive
pairs to probe LLM internals. Each of those requirements is a barrier
in the real-world settings where sentiment analysis is most needed --- new
languages, new domains (legal, medical, financial), and rapid prototyping
pipelines where no labels have been collected yet. A natural question is
whether the \emph{content} that those labels are trying to recover is
already present, in linear form, inside any sufficiently capable
instruction-tuned large language model, and if so whether we can pull it
out without labels.

This paper answers the question affirmatively for valence. We show that a
single unsupervised direction, computed in seconds from nine LLM-generated
paragraphs per class and no human annotations, is a competitive zero-shot
sentiment classifier across five English benchmarks, four non-English
benchmarks, three word-valence lexicons, one continuous VAD regression
task, and two concept-generalisation probes --- seventeen evaluations in
total, all using the \emph{same} 81-story axis.

\paragraph{Our method in one sentence.}
Generate nine short stories for each of nine emotion categories
(\texttt{amusement, inspiration, joy, tenderness, anger, fear, disgust,
sadness, neutral}) using the instruction-tuned LLM itself; extract the
mean-pooled final-layer hidden state of each story; average per class into
nine centroids; take the first principal component of the mean-centered
centroid matrix as the valence axis.

The direction has four properties that together make it a useful
deployment-oriented probe:

\paragraph{1. Unsupervised and data-efficient.}
The axis is extracted with zero labels, zero gradient updates, and zero
human-written contrastive pairs. On SST-2, it matches the accuracy of a
supervised logistic regression trained on $N\!\approx\!100$ labelled
examples on the same hidden features, a $\sim\!10$--$30\times$
data-efficiency gain. A single story per class (9 total samples) already
gives AUC $0.74$ on SST-2; nine stories per class reaches $0.868$.

\paragraph{2. Robust across languages.}
The \emph{English} V-axis, applied zero-shot to held-out non-English
stories, recovers valence at AUC $0.949$ (Spanish), $0.991$ (French),
$0.972$ (German), and $1.000$ (Chinese). Within-language V-axes extracted
from stories in each of the four target languages also transfer above AUC
$0.85$ in every case. Arousal is not preserved in the same way, which
delineates the regime of the method.

\paragraph{3. Robust across model scales and families.}
The identical recipe produces a usable V-axis on Qwen2.5-0.5B, 1.5B, and
7B (all $|r|\!=\!0.48$--$0.49$ on EmoBank V), Pythia-1.4b ($|r|\!=\!0.46$),
and Bloom-560m ($|r|\!=\!0.34$). Across four LLMs the per-story valence
rankings agree at mean cosine $0.961$ and Spearman $0.954$. The method is
not tied to a particular pretraining recipe.

\paragraph{4. Simpler than contrastive alternatives.}
Representation-engineering methods extract concept directions from
operator-specified contrastive pairs (honest/dishonest, happy/sad,
harmful/harmless)
\cite{zou2023representation, rimsky2024steering, arditi2024refusal}. Our
9-centroid PCA \emph{does not require} the operator to supply polarity
labels, and we show it outperforms the mean-difference contrastive axis
(AUC $0.868$ vs $0.753$ on SST-2, $+0.115$) when the same stories are
repartitioned into positive/negative halves.

\paragraph{A scope statement, not a classifier race.}
We are explicit that the method is not competing for state-of-the-art
sentiment accuracy. Direct few-shot prompting of the same LLM reaches
$92.78\%$ on SST-2 ($+15$ percentage points over our median-threshold
probe); a full-rank supervised logistic regression on the \emph{same}
hidden features reaches AUC $0.911$ ($+0.043$). What we offer is the
most data-efficient one-dimensional \emph{summary} of valence that we
could extract under a strict zero-label constraint, together with
evidence that the same summary transfers to many languages and many
tasks. Section~\ref{sec:experiments} quantifies this positioning
explicitly via the counterfactual ablation: removing our axis from the
residual stream changes full-rank LogReg by $0.001$, showing that
valence is distributed across multiple directions, and our axis is the
most useful single-direction extraction under zero supervision.

\paragraph{Where does this direction live?}
A fine-grained layer sweep of Qwen2.5-1.5B reveals a novel ``V-shape''
pattern. The valence direction is \emph{strong} at layer 1
($\mathrm{AUC}\!=\!0.75$), \emph{fades} through the middle layers
($\mathrm{AUC}\!\leq\!0.53$ for layers 2--26), and \emph{suddenly
re-crystallizes} at layer 27 ($\mathrm{AUC}\!=\!0.84$) and 28. The V-axis at
layer 14 is nearly orthogonal (cosine $0.07$) to the V-axis at layer 28,
while L27 and L28 are $95\%$ aligned. The final-layer V-axis is not
inherited from the middle of the network; it is freshly constructed just
before the LM head. This pattern appears in all three causal families we
tested (Qwen, GPT-NeoX, BLOOM) and is absent in bidirectional encoders
(BERT, RoBERTa), whose peak V-axis AUC is at middle layers.

\paragraph{Contributions.}
\begin{itemize}[leftmargin=*]
\item We introduce \ours, a fully unsupervised concept-probing recipe that
requires nine LLM-generated stories per class, no labels, and no gradient
updates. The method produces a linear concept direction in under one minute
on a single V100.
\item We demonstrate that the extracted valence axis transfers to
\emph{nine} languages (EN/ES/FR/DE/ZH plus pairwise cross-lingual
cosines), five LLM families, and seventeen zero-label evaluations with
the same 81-story axis, including SST-2 (AUC $0.868$), IMDB ($0.895$),
Yelp ($0.953$), Twitter SemEval-17 ($0.923$), Rotten Tomatoes ($0.843$),
Hu \& Liu lexicon ($|r|=0.76$), Warriner 3k lexicon ($|r|=0.69$), and
EmoBank continuous valence ($|r|=0.50$).
\item We report a principal-angle analysis showing that PC1 of the nine
centroids captures $55.8\%$ of their variance (vs.\ the next singular value
at $15.2\%$, a $1.9\times$ separation) and lies $89.9^\circ$ from PC2 by
construction; random directions in the same $1\,536$-dim space have mean
angle $88.8^\circ \pm 0.8^\circ$ from PC1, confirming that PC1 is not a
trivial variance-axis artifact.
\item We quantify the method's positioning with a counterfactual ablation:
removing our axis from the residual stream changes supervised full-rank
logistic regression by only $0.001$ AUC, showing that valence is
distributed across many directions. Our contribution is therefore not
``the valence direction'' but the most useful data-efficient
1-D summary of valence extractable under a zero-label constraint.
\item We compare directly to the contrastive mean-difference extraction
of \cite{zou2023representation, rimsky2024steering, arditi2024refusal}: on
the same Qwen-1.5B L28 features, contrastive-pair mean-diff gives SST-2
AUC $0.753$; PCA-of-nine-centroids gives $0.868$ ($+0.115$).
\item We reveal a ``V-shape'' layer-wise emergence pattern: the valence
direction is present at L1, absent through the middle, and re-crystallizes
at the final 1--2 blocks. This pattern replicates across Qwen, Pythia, and
Bloom, and does not appear in bidirectional encoders.
\item We include a scope statement: arousal is not recovered by the same
PCA-of-centroids recipe ($|r|\leq 0.16$ at every layer in every model),
delineating the regime of applicability to valence-like global concepts
aligned with the linear representation hypothesis
\cite{park2024linear}.
\end{itemize}

The rest of the paper: Section~\ref{sec:related} situates \ours{} with
respect to representation engineering, activation steering, and linear
probing; Section~\ref{sec:method} specifies the recipe; Section
\ref{sec:experiments} presents the full zero-label evaluation matrix,
including cross-lingual transfer and the counterfactual ablation; Section
\ref{sec:discussion} discusses when the method works and when it does not.
